\documentclass[12pt,a4paper]{article}

\usepackage[utf8]{inputenc}
\usepackage[T1]{fontenc}
\usepackage[spanish]{babel}
\usepackage{amsmath,amssymb}
\usepackage{geometry}
\usepackage{booktabs}
\usepackage{array}
\usepackage{parskip}
\usepackage{fancyhdr}
\usepackage{listings}
\usepackage{xcolor}
\usepackage{mathtools}

\geometry{margin=2.5cm, headheight=15pt}

\definecolor{codebg}{RGB}{245,245,245}
\definecolor{keywordcolor}{RGB}{0,100,180}
\definecolor{commentcolor}{RGB}{80,110,80}
\lstset{
  language=Python,
  backgroundcolor=\color{codebg},
  frame=single, rulecolor=\color{gray!40},
  basicstyle=\ttfamily\small,
  keywordstyle=\color{keywordcolor}\bfseries,
  commentstyle=\color{commentcolor}\itshape,
  numbers=left, numberstyle=\tiny\color{gray}, numbersep=8pt,
  breaklines=true, showstringspaces=false, tabsize=4,
}

\pagestyle{fancy}
\fancyhf{}
\rhead{Modelado y Simulaci\'on --- 22/04/2026}
\lhead{Punto 2 --- Docente: Omar C\'aceres}
\cfoot{\thepage}

\title{\textbf{Examen --- Modelado y Simulaci\'on}\\[0.3em]
       {\large Punto 2: Interpolaci\'on de Lagrange y Diferencias Finitas}}
\author{Juan Bautista Espino}
\date{22 de abril de 2026}

% ──────────────────────────────────────────────────────────────────────────────
\begin{document}
\maketitle\thispagestyle{fancy}

\section*{Enunciado}

Dada la siguiente tabla de datos discretos:

\begin{center}
\renewcommand{\arraystretch}{1.4}
\begin{tabular}{c|ccccc}
\toprule
$x$ & 0 & 1 & 2 & 3 & 4 \\
\midrule
$y$ & 2 & 6 & $k$ & 12 & 16 \\
\bottomrule
\end{tabular}
\end{center}

\begin{enumerate}
  \item[(a)] Construir el polinomio interpolante de Lagrange con los nodos conocidos y hallar
        el valor de $k$ evaluando $P_n(x)$ en $x_2=2$. Graficar la curva y cada nodo.
  \item[(b)] Usar la tabla generada por el polinomio para calcular $f'(2.5)$ mediante
        diferencias finitas centrales.
\end{enumerate}

% ══════════════════════════════════════════════════════════════════════════════
\section*{Parte (a) --- Polinomio de Lagrange y valor de $k$}
% ══════════════════════════════════════════════════════════════════════════════

\subsection*{Nodos disponibles (excluyendo $x_2=2$)}

Se dispone de cuatro nodos: $(0,2)$, $(1,6)$, $(3,12)$, $(4,16)$.
Con ellos se construye un polinomio de grado $n=3$:

\[
  P_3(x) = \sum_{i=0}^{3} y_i\, L_i(x), \qquad
  L_i(x) = \prod_{\substack{j=0\\j\ne i}}^{3} \frac{x - x_j}{x_i - x_j}
\]

\subsection*{C\'alculo de las bases de Lagrange en $x=2$}

\[
  L_0(2)
  = \frac{(2-1)(2-3)(2-4)}{(0-1)(0-3)(0-4)}
  = \frac{(1)(-1)(-2)}{(-1)(-3)(-4)}
  = \frac{2}{-12}
  = -\dfrac{1}{6}
\]

\[
  L_1(2)
  = \frac{(2-0)(2-3)(2-4)}{(1-0)(1-3)(1-4)}
  = \frac{(2)(-1)(-2)}{(1)(-2)(-3)}
  = \frac{4}{6}
  = \dfrac{2}{3}
\]

\[
  L_2(2)
  = \frac{(2-0)(2-1)(2-4)}{(3-0)(3-1)(3-4)}
  = \frac{(2)(1)(-2)}{(3)(2)(-1)}
  = \frac{-4}{-6}
  = \dfrac{2}{3}
\]

\[
  L_3(2)
  = \frac{(2-0)(2-1)(2-3)}{(4-0)(4-1)(4-3)}
  = \frac{(2)(1)(-1)}{(4)(3)(1)}
  = \frac{-2}{12}
  = -\dfrac{1}{6}
\]

\subsection*{Evaluaci\'on en $x=2$}

\begin{align*}
  P_3(2) &= y_0\,L_0(2) + y_1\,L_1(2) + y_2\,L_2(2) + y_3\,L_3(2) \\
         &= 2\!\left(-\tfrac{1}{6}\right) + 6\!\left(\tfrac{2}{3}\right)
            + 12\!\left(\tfrac{2}{3}\right) + 16\!\left(-\tfrac{1}{6}\right) \\
         &= -\tfrac{2}{6} + \tfrac{12}{3} + \tfrac{24}{3} - \tfrac{16}{6} \\
         &= -\tfrac{1}{3} + 4 + 8 - \tfrac{8}{3} \\
         &= -\tfrac{9}{3} + 12
         = -3 + 12
\end{align*}

\[
  \boxed{k = P_3(2) = 9}
\]

\subsection*{Tabla completa}

\begin{center}
\renewcommand{\arraystretch}{1.4}
\begin{tabular}{c|ccccc}
\toprule
$x$ & 0 & 1 & 2 & 3 & 4 \\
\midrule
$y$ & 2 & 6 & \textbf{9} & 12 & 16 \\
\bottomrule
\end{tabular}
\end{center}

\subsection*{Polinomio interpolante de grado 4}

Con los 5 nodos completos se construye el polinomio de grado 4
(usando diferencias divididas o Lagrange de orden 4):

\[
  P_4(x) = \frac{1}{6}x^4 - x^3 + \frac{29}{6}x + 2
\]

Verificaci\'on: $P_4(0)=2$, $P_4(1)=6$, $P_4(2)=9$,
$P_4(3)=12$, $P_4(4)=16$. $\checkmark$

% ══════════════════════════════════════════════════════════════════════════════
\section*{Parte (b) --- $f'(2.5)$ por Diferencias Finitas Centrales}
% ══════════════════════════════════════════════════════════════════════════════

\subsection*{Tabla densa generada por $P_4(x)$ con $h=0.5$}

\begin{center}
\renewcommand{\arraystretch}{1.3}
\begin{tabular}{cc}
\toprule
$x$ & $P_4(x)$ \\
\midrule
0.0 &  2.000\,000 \\
0.5 &  4.187\,500 \\
1.0 &  6.000\,000 \\
1.5 &  7.562\,500 \\
2.0 &  9.000\,000 \\
2.5 & 10.437\,500 \\
3.0 & 12.000\,000 \\
3.5 & 13.812\,500 \\
4.0 & 16.000\,000 \\
\bottomrule
\end{tabular}
\end{center}

\subsection*{F\'ormula de diferencias finitas centrales}

La f\'ormula de primer orden centrada con paso $h$ es:

\[
  f'(x_0) \approx \frac{f(x_0+h) - f(x_0-h)}{2h}
  + \mathcal{O}(h^2)
\]

\subsection*{C\'alculo en $x_0 = 2.5$, $h = 0.5$}

\[
  f'(2.5) \approx \frac{f(2.5 + 0.5) - f(2.5 - 0.5)}{2 \times 0.5}
  = \frac{f(3.0) - f(2.0)}{1.0}
  = \frac{12.000\,000 - 9.000\,000}{1.0}
\]

\[
  \boxed{f'(2.5) \approx 3.000\,000}
\]

\subsection*{Verificaci\'on con la derivada exacta del polinomio}

Derivando $P_4(x) = \tfrac{1}{6}x^4 - x^3 + \tfrac{29}{6}x + 2$:

\[
  P_4'(x) = \tfrac{2}{3}x^3 - 3x^2 + \tfrac{29}{6}
\]

\[
  P_4'(2.5) = \tfrac{2}{3}(15.625) - 3(6.25) + \tfrac{29}{6}
            = 10.4167 - 18.75 + 4.8333
            \approx 2.9583 \approx 2.958
\]

El error de truncamiento de la f\'ormula centrada es de orden $\mathcal{O}(h^2)$:

\[
  |f'(2.5) - P_4'(2.5)| = |3.000 - 2.958| \approx 0.042
  \sim \frac{h^2}{6}P_4'''(2.5) \quad \checkmark
\]

% ══════════════════════════════════════════════════════════════════════════════
\section*{Ap\'endice: c\'odigo Python}
% ══════════════════════════════════════════════════════════════════════════════

\begin{lstlisting}[caption={Lagrange + diferencias finitas centrales}]
import numpy as np

xs_known = [0, 1, 3, 4];  ys_known = [2, 6, 12, 16]

def lagrange(xs, ys, x):
    result = 0.0
    for i in range(len(xs)):
        Li = 1.0
        for j in range(len(xs)):
            if i != j:
                Li *= (x - xs[j]) / (xs[i] - xs[j])
        result += ys[i] * Li
    return result

# Hallar k
k = lagrange(xs_known, ys_known, 2)   # k = 9

# Tabla con los 5 puntos y h = 0.5
xs_full = [0, 1, 2, 3, 4];  ys_full = [2, 6, 9, 12, 16]
h = 0.5
x_table = [i*h for i in range(9)]
table = {x: lagrange(xs_full, ys_full, x) for x in x_table}

# Diferencias finitas centrales en x=2.5
f_prime = (table[3.0] - table[2.0]) / (2*h)   # = 3.0
\end{lstlisting}

\end{document}
